% WHERE A PARAGRAPH STANDS WHILE A DISPLAY RUNS.
%
% A display breaks the paragraph it interrupts, and the paragraph's own
% level of the nest goes with it: \showlists run inside a display writes the
% display's list and the VERTICAL one under it, with the paragraph's
% \prevdepth and \prevgraf on that vertical level.
%
%   ### display math mode entered at line 11
%   \mathord
%   .\fam1 B
%   ### vertical mode entered at line 0
%   ### current page:
%   ...
%   prevdepth 1.0, prevgraf 1 line
%
% A fresh level is pushed for the paragraph when the display gives it back,
% which is what makes a report of the lines after a display name the line
% the display ended on.
%
% HSTeX kept the paragraph's level standing, so \showlists wrote an extra
% vertical level holding the paragraph's \prevdepth and \prevgraf, and the
% level under it -- the one that really has them -- said `prevdepth
% ignored'.
%
% An \eqno is the same rule the other way round: its list is pushed before
% its mode is changed, so the display's own level stays a DISPLAY math one.
% HSTeX changed the mode first, and trip's \eqno on line 280 turned the
% display it belongs to into `### math mode entered at line 261'.
\catcode`\{=1 \catcode`\}=2 \catcode`\$=3 \catcode`\^=7 \catcode`\_=8
\tracingonline=1 \showboxbreadth=99 \showboxdepth=99
\font\rip=trip \fontdimen22\rip=7pt \rip
\textfont0=\rip \scriptfont0=\rip \scriptscriptfont0=\rip
\textfont1=\rip \scriptfont1=\rip \scriptscriptfont1=\rip
\textfont2=\rip \scriptfont2=\rip \scriptscriptfont2=\rip
\textfont3=\rip \scriptfont3=\rip \scriptscriptfont3=\rip
\hsize=100pt \parindent=5pt \predisplaypenalty=0
\message{[A a display inside a paragraph]}
A A A $$B\showlists$$ C C
\par
\message{[B a display with a group inside it]}
A $${B\showlists}$$
\par
\message{[C inline math in a paragraph, which breaks nothing]}
A $B\showlists$
\par
\message{[D a display with an equation number]}
A $$B\showlists\eqno C\showlists$$
\par
\message{[done]}
\end
